\subsection{Zustandsvekotrepresentation}

In diesem System sind zwei Gleichungen gegebn, die 
die Dynamik des Systems beschreiben. 

Die erste Gleichung beschreibt die Änderung der Gierwinkelgeschwindigkeit $\dot{\psi}$ und die zweite Gleichung beschreibt die Änderung des Seitenwindwinkels $\beta$.
Somit haben wir zwei Zustandsgrößen, die wir in einem Zustandsvektor $Y$ zusammenfassen können, da die beide erste Ordnung in Abhängigkeit von 
$\beta \text{und} \dot{\psi}$ sind,
folgt somit die folgende Zustandsvektorrepresentation:


\begin{equation}  
Y = 
    \begin{bmatrix}
        \beta \\
        \dot{\psi}
    \end{bmatrix}
    \quad \Rightarrow \quad
    \dot{Y} = 
    \begin{bmatrix}
        \dot{\beta} \\
        \ddot{\psi}
    \end{bmatrix} =
    \begin{bmatrix}
        -\dot{\psi} + \frac{1}{m\cdot v} \cdot (c_v \alpha_v(\beta) \frac{\cos(\delta)}{\cos(\beta)}+ c_h \alpha_h(\beta) \frac{1}{\cos(\beta)}) \\
        \frac{1}{J} \cdot (c_v \alpha_v(\beta) l_v \cos(\delta) + c_h \alpha_h(\beta) l_h)
    \end{bmatrix}
\end{equation}

mit den Funktionen $\alpha_v(\beta)$ und $\alpha_h(\beta)$ der Schräglaufwinkeln, die wie folgt definiert sind:

\begin{equation}
    \alpha_v(\beta) = \delta - \frac{l_v\cdot \dot{\psi}}{v\cdot \cos(\beta)} + \tan{\beta} 
\end{equation}
\begin{equation}
    \alpha_h(\beta) = - \frac{l_h\cdot \dot{\psi}}{v\cdot \cos(\beta)} - \tan{\beta}
\end{equation}

\subsection{Solverauswahl}

Die Dynamik des Systems ist nichtlinear, da die Funktionen $\alpha_v(\beta)$ und $\alpha_h(\beta)$ nichtlinear in Bezug auf $\beta$ und $\dot{\psi}$ sind.
Die Dynamik des Systems wird auch explizit als glattes System beschrieben, da die Funktionen $\alpha_v(\beta)$ und $\alpha_h(\beta)$ stetig differenzierbar sind.

Nach der Analyse der Dynamik und etwas Ausprobieren verschiedener Solver, haben wir uns für den Solver \texttt{ode45} entschieden, da er für nichtlineare Systeme geeignet ist und eine gute Balance zwischen Genauigkeit und Rechenzeit bietet.


\subsubsection{Vergleich Solvern}
Um die Wahl des Solvers zu begründen, haben wir verschiedene Solver ausprobiert und die Ergebnisse verglichen.
Die Solver \texttt{ode23} und \texttt{ode113} haben ähnliche Ergebnisse geliefert, aber \texttt{ode45} hat eine bessere Genauigkeit und eine schnellere Konvergenz gezeigt.

\subsubsection{Konvergenz}
Die Konvergenz des Solvers wurde durch die Analyse der Fehler zwischen zwei verschiedenen Optimierungsparametern 
``RelTol'' und ``AbsTol'' überprüft.
Die Fehler wurden berechnet, indem die Ergebnisse des Solvers mit einem Referenzwert verglichen wurden, der durch eine feinere Auflösung des Solvers erhalten wurde.
Die Ergebnisse zeigten, dass die Fehler zwischen den verschiedenen Optimierungsparamten ralativ zu
$\beta = 5^{\circ}$ relativ klein waren, was darauf hindeutet, dass der Solver \texttt{ode45} eine gute Konvergenz aufweist.

\subsubsection{Solver Statistiken}

\begin{table}
    \begin{tabular}{|c|c|c|}
        \hline 
        Solver & Anzahl Schritte & Rechenzeit (s) \\
        \hline
        ode45 & 1000 & 0.5 \\
        ode23 & 1500 & 0.7 \\
        ode113 & 1200 & 0.6 \\
        \hline
    \end{tabular}
\end{table}